\section{判决的效力}
\begin{enumerate}
  \item 既判力。最复杂的一个理论。终局判决一旦确定，该判决对诉讼请求之判断，就成为规范今后当事人之间法律关系的基础。
        当同一事项再度成为问题时，当事人不能对该判断提出争议，不能提出与之矛盾的主张，
        法院也不能作出与该判断相矛盾 \hz 抵触之判断。
        \begin{enumerate}
          \item 生效判决视为真理，禁止再诉的效力。
          \item 我国不承认既判力制度。但考试选项明确出现“既判力”，要选。
          \item 客观范围。什么东西有既判力。
                \begin{enumerate}
                  \item 只及于判决的主文（诉讼请求成立与否），不及于判决的理由。\al{没借过钱；过时效}。
                  \item 部分请求及于全部法律效果。\al{欠 10W，分 10 次起诉，每次要 1W}。
                \end{enumerate}
          \item 主观范围。判决对哪些人具有约束力。
                \begin{enumerate}
                  \item 应当具有相对性。
                  \item 扩张问题。继承、法人合并分立、股东会/董事会决议。
                \end{enumerate}
          \item 时间范围。判决生效时起算（判决生效时的这个点，既判力才形成）。例外。新事实新情况。
        \end{enumerate}
  \item 执行力。原告胜诉的给付之诉的判决可以作为执行根据，具有执行力。\al{有给付内容的生效判决有执行力 \cuo}。
  \item 形成力。判决生效之时即为法律关系变更、消灭之时。
\end{enumerate}
